% !TEX program = xelatex
% ======================================================================
% Banana Quest technical note v4.1 — arXiv submission source
% 编译：含中文摘要，需 XeLaTeX（xeCJK + Noto Sans CJK SC，Overleaf 内置）。
% 若只能用 pdfLaTeX：注释掉 \usepackage{xeCJK}、\setCJKmainfont 两行，
% 并将"中文摘要"环境整段删除（见下方标记 %<PDFFRIENDLY-START/END>）。
% ======================================================================
\documentclass[11pt,a4paper]{article}
\usepackage{xeCJK}
\setCJKmainfont{Noto Sans CJK SC}
\usepackage[margin=2.5cm]{geometry}
\usepackage{booktabs}
\usepackage{array}
\usepackage{xcolor}
\usepackage{microtype}
\usepackage[colorlinks=true,linkcolor=black,urlcolor=blue,citecolor=black]{hyperref}

\definecolor{linkblue}{HTML}{1A5FB4}
\newcommand{\ds}[1]{\textcolor{linkblue}{\url{#1}}}
\newcolumntype{L}[1]{>{\raggedright\arraybackslash}p{#1}}

\title{Dynamics on aggregated connectomes in the browser:\\
normalization, structural bias, and honest tuning lessons\\
across two connectomes}
\author{Banana Quest project (\texttt{fruitfly/game/})}
\date{Technical note v4.1, 2026-09-27}

\begin{document}
\maketitle

\noindent\small
\textbf{Technical note v4.1, 2026-09-27.} Code: \texttt{fruitfly/game/} --- v1
measurements @ commit \texttt{d92807b}, v2 probes @ \texttt{73da541}, v3--v4
scripts \& experiments @ \texttt{63c54ca} / \texttt{709b188}, this revision @
\texttt{80a4768} (MIT). Data: FlyWire FAFB v783 (CC BY-NC-SA 4.0) and
MaleCNS v1.0 (CC BY 4.0). v4.1 closes the third review round (``minor
revisions, no new experiments required''): histamine-account correction,
reverse-experiment precision, two final frozen reproduction entries, and
editorial cleanup. Point-by-point responses: \texttt{docs/REVIEW-RESPONSE.md}.
Repository: \ds{https://github.com/professorwang/flybrain-banana-quest};
demo: \ds{https://professorwang.github.io/flybrain-banana-quest/}.
\normalsize

%<PDFFRIENDLY-START> —— pdfLaTeX 备选：删除本段中文摘要即可
\section*{中文摘要}
本文记录浏览器游戏《电子果蝇·香蕉大作战》（Banana Quest）在真实连接组上跑动力学时实测到的方法学问题及其跨数据集复测结果，供连接组仿真社区参考。游戏用 LIF 脉冲网络在两个真实连接组（FlyWire FAFB v783：139,255 神经元 / 2,698,236 条聚合连接；MaleCNS v1.0：176,422 神经元 / 6,287,749 条 $\ge$5 突触连接）上驱动一只 2D 果蝇觅食——连接组放电参与运动读出与进食门控，同时游戏包含显式运动规则（\S1.1 如实列出）。发现一：参考实现沿用的全局 max 权重归一化（增益 0.15）在\textbf{参考增益下}传播失败——重尾突触计数分布把典型权重压到阈值的约 1/2000，下游膜电压在 100 ticks 内观察到的最大值仍不过阈（FlyWire 0.780、MaleCNS 0.977），中枢放电为零；增益放大 10 倍后两数据集分别有 16 / 131 次中枢放电，因此结论收窄为“参考增益下失败”而非“必然失败”。可行替代是本项目采用的 postsynaptic L1 归一化（每个突触后神经元总入权重归一化为常数；与 Shiu et al.\ 2024 的每突触固定系数方案不同），工作点为数据集相关（FlyWire 3.0 / MaleCNS 4.0）。发现二：FlyWire 下行池按胞体 x 切半后存在结构性偏置（任一侧刺激右池都更活跃，有效差仅约 1.3\%），转向读出采用本项目选用的启发式高通读出（慢速 EMA 基线的偏差）；该偏置在所测 MaleCNS 条件下未复现（侧向对比方向正确，差 9.49 个百分点，算式 $91/157 - 111/229 \approx 0.09490$），“分组伪影”解释目前仅为待验证假说。发现三（方法论）：小种子扫参会过拟合——实测冠军配置在 10 个样本外种子上未保持优势（首吃时间只统计成功 episode）。符号实验闭环：仅把 MaleCNS 的谷氨酸符号从抑制翻为兴奋，ti=4 时视叶出现高活动（VIS\_OL 553,863 次放电）；反向把 FlyWire 的谷氨酸翻为抑制，ti=4 视叶高活动消失（VIS\_ME 405,788 $\to$ 0）——在 ti=4、I=1、200 tick 及所比较的 GLUT 正/负两种配置下，谷氨酸符号赋值既是高视叶活动表型的充分条件也是必要条件（ti=3 结果照实报告但不纳入该结论范围）；两数据包符号表不同（GLUT 兴奋 vs 抑制；MaleCNS 明确将 8,021 个 histamine 标记节点设为抑制，而 FlyWire 的六类递质预测不含 histamine——不能从缺少该预测类别直接推断真实组胺能神经元在 FlyWire 包中的赋号），是跨数据包动力学比较必须显式声明的混杂。所有数字均可由文中命令复现（注明历史记录者除外）。
%<PDFFRIENDLY-END>

\begin{abstract}
\noindent Banana Quest is a zero-dependency browser game in which a leaky
integrate-and-fire (LIF) network running on a real connectome helps drive a 2D
foraging fly. Connectome activity participates in the motor readout and
feeding gating; the game additionally contains explicit motor rules, which we
list openly (\S1.1). We report three methodological findings and their
cross-dataset replication on FlyWire FAFB v783 and MaleCNS v1.0. (1)
Global-max weight normalization \textbf{at the reference gain (0.15)} fails to
propagate activity on aggregated connectomes with heavy-tailed synapse-count
distributions --- the maximum membrane voltages observed within 100 ticks stay
below threshold (0.780 / 0.977), with zero central spikes on both datasets; at
$10\times$ gain a small number of central spikes appears (16 / 131), with
unknown stability, so we scope the claim to the reference gain rather than
``global-max never propagates''. A working alternative is postsynaptic L1
normalization (adopted in this project; distinct from Shiu et al.'s fixed
per-synapse coefficient scheme), with a dataset-dependent working point (3.0 /
4.0). (2) A coarse bilateral descending readout on FlyWire carried a
structural bias larger than the stimulus-driven signal ($\sim$1.3\%); it did
not replicate under the tested MaleCNS conditions (correctly signed contrast,
9.49 percentage points, $91/157 - 111/229 \approx 0.09490$), so the
``pool-definition artifact'' explanation remains a hypothesis --- the decisive
test is re-pooling on the same dataset. (3) Small-seed sweeps of high-variance
behavioral metrics can overfit: concretely, our 5-seed champion configuration
did not retain its advantage on 10 out-of-sample seeds (first-eat statistics
count successful episodes only). Sign experiments now close both directions:
flipping only the glutamate sign from inhibitory to excitatory in MaleCNS
produces high visual-lobe activity at ti=4 (553,863 VIS\_OL spikes), and the
reverse flip in FlyWire abolishes it (VIS\_ME 405,788 $\to$ 0) ---
\textbf{at ti=4, I=1, 200 ticks, and across the two GLUT sign configurations
compared}, the glutamate sign assignment is both sufficient and necessary for
the high-activity phenotype (ti=3 results are reported but excluded from this
conclusion). The two packages also differ in their sign tables: GLUT
excitatory vs inhibitory, and MaleCNS explicitly marks 8,021 histamine-labeled
nodes inhibitory while FlyWire's six-class transmitter prediction contains no
histamine class at all --- the absence of a predicted class cannot be used to
infer how real histaminergic neurons are signed in the FlyWire package. These
are confounds any cross-package dynamical comparison must declare. All numbers
are reproducible with the commands given, except where marked as historical
records.
\end{abstract}

\section{Background}

Banana Quest (电子果蝇·香蕉大作战) is an open-source, zero-build browser game.
A virtual fly forages for bananas in a 2D arena. Olfactory input is injected
as current into left/right pools of olfactory receptor neurons; spikes
propagate through the connectome to descending neurons, whose pooled
left/right firing-rate difference steers the fly, and whose total rate, scaled
by a hunger variable, sets its speed. Feeding is gated by gustatory
stimulation driving a proboscis motor-neuron readout. \textbf{Connectome
spiking participates in steering and in feeding gating --- but it is not the
only thing moving the fly: the game also contains explicit motor rules,
disclosed in \S1.1.}

The FlyWire connectome ships as a compact binary (12.4 MB gzipped) packaged by
the community project snedea/flybrain (MIT): FAFB v783 edges aggregated per
neuron pair with signed weights, plus a coarse 63-group functional annotation.
Their reference simulator runs the full network in a Web Worker at
$\sim$10 ticks/s using a neuropil-gated dormant-group optimization; our
simulation core (\texttt{src/sim-core.js}) is adapted from it and runs
unchanged in both the browser worker and Node.js for testing. The second
dataset is MaleCNS v1.0 (neuPrint \texttt{male-cns:v1.0}; 176,422 neurons,
6,287,749 connections at $\ge$5 synapses; the full male central nervous system
including the VNC; CC BY 4.0), selected at runtime with
\texttt{?dataset=malecns}. We take the FLYB v1 binary bundled by the
fly-brain-minecraft project (MIT) and convert it with our own
\texttt{tools/flyb\_to\_bin.py} into the same sim-core format --- 26
functional groups with the same region encoding, and pools filtered on
MaleCNS's native annotations (somaSide, superclass, type). Every probe below
runs through the identical \texttt{sim-core.js} on both datasets.

\subsection{Explicit game-layer motor rules}

Earlier versions of this note said ``no hand-crafted behavior layer'' --- that
was overstated. The brain participates in behavior, but the following explicit
rules also move or stop the fly, and reviewers rightly asked for them to be
listed:

\begin{itemize}
\item \textbf{Standby random walk}: when the descending pools' total firing
rate falls below a threshold, the fly performs a slow random walk driven by
game-layer noise, explicitly labeled ``待机噪声（非脑驱动）'' in the UI.
\item \textbf{Approach deceleration and stopping}: speed is additionally
scaled down when turning sharply and when close to the banana (a
turn-and-slow kinematic rule), and the fly stops inside the eating radius
while the feeding gate is evaluated.
\item \textbf{Scripted escape maneuver}: after a startle, the escape
\emph{trigger} is the brain's response (descending-rate surge or DNp01
firing), but the escape \emph{maneuver} itself --- a random sharp turn plus a
fixed-duration speed boost --- is scripted.
\item \textbf{Wall bounce}: heading reflects off arena walls with added noise.
\item \textbf{Hunger wiring}: hunger is a scripted clock that multiplies
olfactory gain and injects current into a hunger-drive group (an acknowledged
artificial wiring).
\item \textbf{Sensor mapping and pool selection}: the odor$\to$current
mapping, the left/right pool splits, and all readout pools are hand-designed
(acknowledged throughout; MaleCNS stimulus pools are additionally asymmetric,
\S5.5).
\end{itemize}

What the brain does do: it gates eating (the fly cannot eat unless the
proboscis readout fires), it modulates steering and speed through the
descending pools, and it decides whether a startle produces escape. What it
does not do: fine motor coordination of walking, which is VNC-driven in real
flies and kinematic rendering here.

\section{Finding 1 --- Global-max normalization at the reference gain does not propagate on aggregated connectomes}

\subsection{Measurement}

The reference implementation normalizes edge weights globally:
$w \leftarrow w / \max|w| \times 0.15$. On both aggregated datasets the raw
synapse-count weights are extremely heavy-tailed:

\begin{center}
\begin{tabular}{lcccc}
\toprule
dataset & max & median & p90 & p99 \\
\midrule
FlyWire v783 & 2405 & 8 & 23 & 77--83 (sample-dependent) \\
MaleCNS v1.0 & 2591 & 9 & 28 & 93 \\
\bottomrule
\end{tabular}
\end{center}

Global-max normalization therefore maps the \emph{median} edge to
$\approx 0.15 \times 8/2405 \approx 5\times10^{-4}$ --- roughly 1/2000 of the
firing threshold (1.0). Empirically, driving the left olfactory pool at
sustained intensity 0.5 for 100 ticks produced, on \textbf{FlyWire}, 18,525
spikes \textbf{all in sensory groups; central-brain spikes: 0}, with the
downstream OLF\_PN group's maximum membrane voltage observed within those 100
ticks reaching only \textbf{0.780} (0.688 within the first 50 ticks of the
same run). On \textbf{MaleCNS} the same protocol produced 21,867 spikes, again
all sensory, central and motor 0, with the 100-tick observed maximum at
\textbf{0.977}. Signal dies in the antennal lobe at the reference gain on both
datasets.

\subsection{Postsynaptic L1 normalization (adopted in this project)}

We replaced global scaling with \textbf{postsynaptic L1 normalization}: each
postsynaptic neuron's total incoming $|w|$ is scaled to a constant
\texttt{targetInput}. \emph{Attribution note: earlier versions of this note
called this ``Shiu et al.\ (2024)-style'' --- a misattribution. Shiu et al.'s
published code (\texttt{philshiu/Drosophila\_brain\_model}, \texttt{model.py})
uses a fixed per-synapse coefficient (\texttt{w\_syn = 0.275 $\cdot$ mV}), not
postsynaptic total-input normalization. The scheme adopted here is our own
engineering choice for this project; the two differ fundamentally (per-neuron
budget vs global coefficient).} Sweeping \texttt{targetInput} (200-tick probes
driving the same ORN pool):

\begin{center}
\begin{tabular}{lll}
\toprule
targetInput & FlyWire (200 ticks) & MaleCNS (200 ticks) \\
\midrule
1.0 & MB\_KC 2,544, GNG\_DESC 0 & MB\_KC 148, DN 0 \\
2.0 & MB\_KC 57,446, GNG\_DESC 10 & MB\_KC 37,237, DN 0 \\
3.0 & GNG\_DESC 3,810; stable over 500 ticks of maximal bilateral input & DN 11; no visual ignition \\
4.0 & whole-brain ignition (VIS\_ME 405,788 spikes/200 ticks, zero visual input) & DN 75; VIS\_OL 2 (no ignition) \\
\bottomrule
\end{tabular}
\end{center}

The working point is dataset-dependent: we ship 3.0 for FlyWire and 4.0 for
MaleCNS (\S5.3).

\subsection{Gain dependence (review replication)}

An external reviewer noted that the failure claim depends on the gain
constant. Repeating the \S2.1 protocol with the global-max coefficient scaled
$\times 10$ ($0.15 \to 1.5$): \textbf{a small number of central spikes appears
on both datasets --- 16 on FlyWire, 131 on MaleCNS (plus 13 motor)}
(\texttt{node game/tools/probe\_gain.mjs}; archive
\texttt{docs/results/probe\_gain.txt}). We state this narrowly: at reference
gain 0.15, propagation fails on both datasets; at $\times 10$, a small number
of central spikes appears, and the stability or usefulness of the higher-gain
regime is unknown (not characterized here). Accordingly, Finding 1 is scoped
to the reference gain throughout.

\subsection{Discussion}

Aggregation is the culprit multiplier at the reference gain: collapsing all
synaptic contacts between a neuron pair into one scalar produces a few huge
weights (a 2405-contact edge) that set the global scale for everyone else, and
a threshold model then divides the world into ``edges that matter'' and
``edges that don't'' with nothing in between. Any demo that ``just runs
dynamics'' on an aggregated connectome should publish its normalization, its
gain, and a propagation test; a silent network at reference gain is the
default outcome, not a bug in your stimulus.

\section{Finding 2 --- Structural bias in coarse descending readouts (FlyWire measurement)}

\subsection{Measurement}

Our FlyWire steering readout splits GNG\_DESC (3,581 neurons) into left/right
pools by soma x-coordinate median (1,791 / 1,790 neurons). Stimulating either
antennal pool alone (I=1.0, 200 ticks, targetInput 3.0) gives:

\begin{center}
\begin{tabular}{lccc}
\toprule
stimulus & desc\_left & desc\_right & L/(L+R) \\
\midrule
left antenna pool & 1,704 & 2,106 & 0.447 \\
right antenna pool & 1,633 & 2,133 & 0.434 \\
\bottomrule
\end{tabular}
\end{center}

The right pool fires more under \textbf{either} unilateral stimulus (bias
$\approx +24\%$ over the left pool), while the stimulus-dependent contrast
between the two rows is only $\sim$1.3\%. First descending spikes appear 9
ticks ($\sim$0.9 s at 10 Hz) after stimulus onset. Feeding the raw rates into
a steering law $\mathrm{turn} = k\cdot(R-L)/(R+L+\varepsilon)$ would make the
fly veer right forever regardless of where the banana is.

\subsection{Heuristic high-pass readout}

We steer on a \textbf{heuristic high-pass readout adopted in this project}:
$\mathrm{turn} = k\cdot(\hat{R}-\hat{L})/(\hat{R}+\hat{L}+\varepsilon)$ with
$\hat{R} = \max(0, R - \mathrm{EMA}_\tau(R))$, $\tau = 5$ s. We make no claim
that this removes the bias or recovers a ``true'' lateral signal; it is a
heuristic chosen for responsiveness, and the only ablation we offer is
behavioral: longer baselines ($\tau = 20$--$40$ s) perform worse in our sweeps
(eat rates fall to 40--80\% in 7 of 8 swept configurations, and at $\tau = 40$
s path tortuosity rises to $\sim$17--23), which we interpret as persistent
wrong-sign deviations getting ``locked in'' by the long memory. Whether any
baseline-free readout can recover steering contrast on this dataset is an open
experiment (\S6).

\subsection{Interpretation --- hypothesis, not conclusion}

Where does the bias come from? Candidate contributors: the fly's largely
bilateral olfactory pathway, our anatomically naive x-median split, the
``GNG\_DESC'' group definition mixing descending and ascending neurons, and
single-specimen hemispheric asymmetries in proofreading or annotation. v2
leaned toward ``real structure''; v2's cross-dataset result (\S5.4) instead
suggests a pool-definition artifact. In v3 we explicitly downgrade all such
attribution to \textbf{hypothesis}: the measurement itself stands, the
explanation is undecided, and the decisive experiment is
re-pooling/relabeling on the same FlyWire data (\S6). Practical rule for demo
builders, unchanged: \textbf{calibrate readout pools against a slow baseline
and report the bias} --- if your bilateral pools respond equally to both
sides, your ``steering'' is cosmetic.

\section{Finding 3 --- Seed overfitting: a tuning-methodology lesson}

\subsection{The sweep winner that wasn't}

Because steering is weak and fluctuation-dominated, time-to-first-banana is a
high-variance metric (in our runs, 12--166 s across seeds for identical
parameters; the game injects constant current, so the fluctuations arise from
network dynamics and game-layer randomness, not from any Poisson
implementation --- v2's ``Poisson noise'' wording was imprecise). We swept
game-layer readout constants (\texttt{turnGain} $\times$ \texttt{smellSigma}
$\times$ \texttt{baseSpeed}, then \texttt{turnGain} $\times$ \texttt{emaTau}
$\times$ \texttt{baseSpeed}; 12 configs $\times$ 5 LCG-seeded runs $\times$
180 simulated seconds each, brain reset between episodes, full tables in
\texttt{game/README.md}). The qualitative trends are modest and
configuration-dependent: in our tables, higher \texttt{turnGain} is associated
with worse outcomes in most rows, but there are counterexamples --- e.g.\ at
fixed \texttt{smellSigma}=250 and \texttt{baseSpeed}=12, raising
\texttt{turnGain} from 2.6 to 4 \emph{improves} mean score from 1.80 to 2.20
--- and enlarging \texttt{emaTau} degrades performance (\S3.2). All first-eat
statistics below count \textbf{successful episodes only} (a run with no eat
contributes to the eat-rate but not to first-eat time). The apparent champion,
\texttt{turnGain} = 1.8, scored 100\% eat rate, 26.4 s median, 3.00 mean
score across the 5 sweep seeds.

On 10 fresh seeds it collapsed: 8/10 ate, median 57.7 s, mean 73.7 s, mean
score 1.30 --- \textbf{worse than the status-quo default on the same 10 seeds}
(10/10, 49.6 s median, 2.70 score). The 5-seed ``win'' was seed luck. A
configuration-reproducibility caveat adds a second layer: the champion's 3.00
was measured when \texttt{bananaMaxDist} was still unbounded; under the
current default (\texttt{bananaMaxDist} = 250) the same configuration scores
2.20, and explicitly restoring the historical unbounded value recovers 3.00
(both numbers match the reviewer's independent measurements; commands in \S7).

\subsection{What actually helped: level design}

Capping banana spawn distance at 350 px excluded the far corners of the arena
(max distance from the center-spawned fly is $\approx$362 px), yet produced
metrics identical to the default in our seed set --- no measurable effect.
Capping spawns to a 150--250 px ring, however, moved the median
time-to-first-banana from 49.6 s to \textbf{27.3 s} (10/10 eaten, tortuosity
8.47 --- still clearly wandering search rather than a $\approx$1 straight-line
beeline). The ``fast search'' family (higher speed gain / max speed, rewriting
the rate$\to$speed readout mapping) tells a mixed story and was rejected:
fast-search alone improved the median (38.6 s $<$ 49.6 s) but \emph{reduced}
the eat rate to 9/10 and the mean score to 1.80; ring+fast combined was worse
than ring-only on the median (42.6 s $>$ 27.3 s) while scoring 3.10 --- a
score gain we attribute to faster wandering covering the ring sooner, not to
any readout improvement, and we did not adopt it.

\subsection{Honest-tuning principle}

The final shipped defaults change \emph{only} the level-design constant
(\texttt{bananaMaxDist} = 250). Every constant in the
sensory$\to$brain$\to$readout mapping is unchanged, so the demo's core claim
--- ``the brain's readout participates in the decision'' --- carries no hidden
tuning water. Two methodological lessons generalize beyond games. First:
\textbf{behavioral metrics with high stochastic variance must be validated
out-of-sample; a small-seed sweep does not estimate the parameter effect, it
estimates the seed set.} Second, disclosed in v3: \textbf{the seeds (101--110)
used to select the level-design constant were also used to evaluate it ---
they are no longer an independent test set.} We kept both tables in the repo
README precisely so the selection is auditable.

\section{Cross-dataset replication on MaleCNS}

Everything in \S2--\S4 was first measured on FlyWire. We then re-ran the same
probe battery on MaleCNS v1.0 through the identical simulation core. Full
tables and commands are in \texttt{docs/malecns-probes.md} and
\texttt{docs/results/}.

\subsection{Preprocessing caveats}

\textbf{Both datasets are $\ge$5-synapse filtered.} v2 claimed the FlyWire
package ``keeps aggregated connections of any strength'' --- that was wrong.
Unsigned aggregation of the FlyWire source CSVs yields a minimum connection
weight of 5.0 with zero neuron pairs below 5 (2,700,513 unsigned pairs; the
2,277-pair difference to the binary's 2,698,236 edges is pairs whose signed
weights cancel to exactly zero and are dropped). The MaleCNS bundle is
likewise thresholded at $\ge$5 synapses (6,287,749 connections) and
additionally drops autapses.

\textbf{The neurotransmitter sign rules differ between packages --- a major
confound.}

\begin{center}
\begin{tabular}{lll}
\toprule
transmitter & FlyWire package (snedea \texttt{NT\_SIGN}) & MaleCNS package (fly-brain-minecraft) \\
\midrule
acetylcholine & $+1$ excitatory & $+1$ excitatory \\
glutamate & $\mathbf{+1}$ \textbf{excitatory} & $\mathbf{-1}$ \textbf{inhibitory} \\
histamine & \textbf{not a predicted class in this dataset} & \textbf{$-1$ inhibitory (8,021 labeled nodes)} \\
GABA & $-1$ inhibitory & $-1$ inhibitory \\
monoamines (DA/OA/5-HT) & $+1$ excitatory & $+1$ excitatory \\
\bottomrule
\end{tabular}
\end{center}

Correction history on histamine (kept for transparency): v3 incorrectly
claimed both packages treat histamine as inhibitory; v4 over-corrected by
adding ``no entry $\to$ default $+1$, photoreceptors directly affected''. Both
were wrong in different ways. The facts, verified by scanning every
\texttt{nt\_type} label in the FlyWire source \texttt{connections.csv.gz}
(3,869,878 rows): the label set is exactly six classes --- ACH 2,258,155; GABA
865,318; GLUT 654,183; DA 37,705; SER 37,450; OCT 17,067 --- \textbf{with no
histamine and no empty labels}. FlyWire's six-class transmitter prediction
simply does not include histamine, and neurons (including photoreceptors) are
labeled with one of these six classes; \textbf{the absence of a predicted
class cannot be used to infer how real histaminergic neurons are signed in the
FlyWire package}. The MaleCNS package, by contrast, explicitly assigns $-1$ to
8,021 histamine-labeled nodes (that count is the total of histamine-labeled
nodes, not a photoreceptor count). Glutamatergic neurons are abundant in both
datasets (MaleCNS: 29,763 neurons, $\sim$17\%). Cross-dataset dynamical
differences below must be read against these sign-table differences. Other
differences remain: sex (female/male), VNC (absent/present), individuals, and
packaging pipelines.

\subsection{Finding 1 --- replicated at reference gain, and gain-bounded}

Covered in \S2.1 and \S2.3: global-max normalization fails at the reference
gain on both datasets (100-tick observed maxima 0.780 / 0.977, zero central
spikes), and a small number of central spikes appears at $10\times$ gain (16 /
131 central spikes). The cross-dataset statement that survives review is
therefore narrow but robust: \textbf{at the reference gain of the reference
implementation, global-max normalization does not propagate on either
aggregated connectome.}

\subsection{targetInput --- structure replicates, the working point drifts; sign rules, not filtering, drive the ignition difference}

The per-neuron sweep (\S2.2 table) reproduces the same layered reachability on
MaleCNS --- 1.0 dies at the mushroom-body entrance, 2.0 ignites the mushroom
body but not descending neurons, 3.0 first reaches them (DN 11), 4.0 makes
them usable (DN 75). Two quantitative differences from FlyWire:

First, the olfactory$\to$descending-neuron drive is \textbf{about two orders
of magnitude weaker} (11 vs 3,810 GNG\_DESC spikes per 200 ticks at ti=3.0).
Candidate contributors: MaleCNS's \texttt{descending\_neuron} superclass is a
pure brain$\to$VNC command population several synapses from the olfactory
entry, whereas FlyWire's ``GNG\_DESC'' group mixes in directly sensory-driven
neurons; and the sign-rule difference changes net excitation everywhere.

Second --- and here v2's attribution must be retracted --- \textbf{MaleCNS
shows no FlyWire-style high visual-lobe activity at ti=4.0} (VIS\_OL: 2
spikes, versus VIS\_ME 405,788 spikes/200 ticks on FlyWire). v2 attributed
this to MaleCNS's $\ge$5-synapse filtering; that attribution was void (\S5.1:
both datasets are filtered). The replacement hypothesis was the sign-table
difference, and we tested it in both directions:

\textbf{Forward experiment (sufficiency, MaleCNS).} We regenerated the MaleCNS
binary with only the glutamate sign flipped from $-1$ to $+1$
(\texttt{python game/tools/flyb\_to\_bin.py {-}{-}glut-excitatory}, 29,763
neurons flipped; all other parameters identical), and ran the ti=3/4 probes:

\begin{center}
\begin{tabular}{lccc}
\toprule
variant & total spikes (200 ticks) & DN & VIS\_OL \\
\midrule
MaleCNS standard (GLUT inhibitory), ti=3.0 & 235,638 & 11 & 0 \\
MaleCNS standard, ti=4.0 & 331,336 & 75 & 2 \\
MaleCNS GLUT-excitatory, ti=3.0 & 492,145 & 1,025 & 41 \\
\textbf{MaleCNS GLUT-excitatory, ti=4.0} & \textbf{1,797,262} & \textbf{19,004} & \textbf{553,863} \\
\bottomrule
\end{tabular}
\end{center}

\textbf{Reverse experiment (necessity, FlyWire).} We rebuilt the FlyWire
binary from the source connection table with only the glutamate sign flipped
from $+1$ to $-1$ --- i.e.\ the source rows are re-aggregated and
postsynaptic-L1-normalized identically, differing only in the sign rule
(\texttt{python game/tools/build\_flywire\_variant.py}; equivalence
self-check: a standard-sign rebuild reproduces the shipped
\texttt{connectome.bin.gz} edge-for-edge, 2,698,236 edges, zero mismatches),
and ran the same probes:

\begin{center}
\begin{tabular}{lccc}
\toprule
variant & total spikes (200 ticks) & GNG\_DESC & VIS\_ME \\
\midrule
FlyWire standard (GLUT excitatory), ti=3.0 & 447,219 & 3,810 & 2,860 \\
FlyWire standard, ti=4.0 & 1,265,998 & 25,730 & 405,788 \\
FlyWire GLUT-inhibitory, ti=3.0 & 222,036 & 30 & \textbf{0} \\
FlyWire GLUT-inhibitory, ti=4.0 & 291,831 & 113 & \textbf{0} \\
\bottomrule
\end{tabular}
\end{center}

\emph{Scope note.} The intervention is a single rule change, not a single-edge
change: it alters the edge set by 517 edges (2,698,753 vs 2,698,236, from
changed cancellation), the weights of 77,650 shared edges, and 23,905
postsynaptic L1 denominators. These are consequences of the same rule
intervention, not an additional confound --- but ``the graph'' here therefore
means the fixed source connection table re-aggregated under each sign rule,
not one frozen weighted graph.

What these two experiments establish, stated as narrowly as the data allow:
\textbf{at ti=4, I=1, 200 ticks, and across the two GLUT sign configurations
compared, flipping only the glutamate sign is sufficient to cause high
visual-lobe activity on the MaleCNS-derived graph, and flipping it back
abolishes that activity on the FlyWire-derived graph (405,788 $\to$ 0) ---
the glutamate sign assignment is both sufficient and necessary for the
high-activity regime under exactly these conditions.} The ti=3 rows are
reported for completeness and are \emph{not} part of this conclusion. This
directly supports sign assignment as a major contributing factor to the
cross-package dynamical difference; it does not prove sign tables explain
\emph{all} cross-dataset differences (drive asymmetry, anatomy, and filtering
still differ). Excitability is partly a packaging choice: Shiu et al.\ assign
glutamate to the inhibitory category (GluCl), the FlyWire browser package to
the excitatory one, and the resulting networks sit in measurably different
dynamical regimes. (\texttt{node game/tools/probe\_signflip.mjs}; archives
\texttt{docs/results/probe\_signflip.txt},
\texttt{docs/results/signflip\_flywire.txt}.)

The practical advice to the community sharpens accordingly: \textbf{run a
per-dataset propagation test, publish the working point and the sign table ---
there is no universal constant, and there is no universal sign convention
either.}

\subsection{Finding 2 --- not replicated under the tested MaleCNS conditions}

With MaleCNS's \emph{native} \texttt{somaSide} pool split and adequate drive
(ti=4.0), the lateral contrast is \textbf{correctly signed}:

\begin{center}
\begin{tabular}{llccc}
\toprule
drive & stimulus & desc\_left & desc\_right & L/(L+R) \\
\midrule
ti=3.0, 200 ticks & left antenna pool & 5 & 6 & 0.455 \\
ti=3.0, 200 ticks & right antenna pool & 12 & 15 & 0.444 \\
ti=4.0, 400 ticks & left antenna pool & 91 & 66 & \textbf{0.580} \\
ti=4.0, 400 ticks & right antenna pool & 111 & 118 & \textbf{0.485} \\
\bottomrule
\end{tabular}
\end{center}

The left--right difference between the two ti=4 rows is
\textbf{$91/157 - 111/229 \approx 0.09490$ (9.49 percentage points)}. At
ti=3.0 the DN counts are too small (5--15 spikes per pool in 200 ticks;
10--34 in 400 ticks) for any bias statistic.

What can and cannot be concluded. Under the tested MaleCNS conditions, the
wrong-direction bias does not appear. But the two measurements differ in
dataset, individual, sex, VNC presence, \textbf{sign tables}, pool
definitions, and working point all at once --- so ``pool-definition artifact''
(v2's preferred reading) remains \textbf{a hypothesis, not a conclusion}. The
decisive experiment changes only the pooling on a fixed dataset: re-split
FlyWire's descending population by alternative rules (e.g.\ hemispheric
relabeling, ascending-excluded re-pooling) and re-measure (\S6, next
experiments). One related observation, kept strictly within MaleCNS: its
ascending neurons (AN group) fired 0 times in 200 ticks under olfactory
stimulation (versus 710 under bristle stimulation). Whether FlyWire's mixed
GNG\_DESC group derives any of its olfactory responsiveness from its own
ascending component \textbf{has not been verified} --- we removed the v3
sentence that claimed otherwise.

\textbf{Statistical disclosures.} (i) The MaleCNS olfactory stimulus pools are
asymmetric: 411 unknown-side ORNs were assigned to the left/right pools by
index parity, giving 1,090 (left) vs 1,549 (right) neurons --- at equal
per-neuron current, the right pool receives $\approx$42\% more total input.
(ii) 10 midline descending neurons were assigned to the pools by parity.
(iii) Each probe condition is a single deterministic trajectory; spike counts
within a run are not independent samples, and ti=4.0 does not by itself make
the comparison statistically adequate. (iv) The heuristic high-pass readout
of \S3.2 remains the readout used in the game on both datasets.

\subsection{Biological sanity checks (small, but free)}

\begin{itemize}
\item \textbf{DNp01 (giant fiber) does not respond to bristle stimulation}: 0
spikes in 20 ticks of strong mechanosensory-bristle drive --- consistent with
its known biology as a visually (looming-) triggered escape neuron. The game's
escape path therefore keys on the descending-rate surge (bristle drives
ascending/central activity strongly, 710 AN spikes in the same probe), not on
DNp01.
\item \textbf{The feeding chain works on MaleCNS}: gustatory pool at I=1.2
drives the MN9-class proboscis motor readout (16 cb\_motor neurons) to 3
spikes in 30 ticks (first at tick 22) --- at the game's feeding-gate
threshold, and the headless closed loop on MaleCNS eats bananas for real
(score 1--3 per 180 s across seeds, first eats 49--172 s; steering is
legitimately ``hard mode'': DN rates of $\sim$2--8 spikes/s at distance,
rising to 34--42/s at close range).
\end{itemize}

\section{Limitations and open questions}

Limitations. Point LIF neurons without conductances, delays, neuromodulation,
or plasticity; no gap junctions; transmitter signs from prediction pipelines
with package-specific sign rules that materially change dynamics (\S5.1,
\S5.3); two individuals (one female FAFB, one male MaleCNS) with different
packaging pipelines; FlyWire has no ventral nerve cord, and even on MaleCNS we
do not simulate real gait; \textbf{the game contains explicit motor rules
(\S1.1)}; per-dataset stimulus gains are acknowledged level design; probe
conditions are single deterministic trajectories, and the level-design seed
set doubles as its evaluation set (\S4.3). None of this claims biological
fidelity, let alone ``uploading''.

Next experiments (ordered):
\begin{enumerate}
\item \textbf{Decisive test for the Finding 2 attribution}: on the \emph{same}
FlyWire data, re-pool the descending population under alternative rules
(hemispheric relabeling, ascending-excluded, native-side annotations from
FlyWire's own classification) and re-measure the unilateral-stimulus contrast.
If the wrong-direction bias vanishes under alternative pooling, the artifact
reading is confirmed on that dataset.
\item \textbf{BANC replication} of the global-max reference-gain failure
(one-command dataset swap), plus a record of each packaging's sign table as a
mandatory companion to any dynamical comparison.
\item \textbf{Named-DN readouts}: which specific descending types (e.g.\ the
turn/halt DNs characterized in the literature) restore steering contrast
without any baseline subtraction, on either dataset?
\item \textbf{Gain stability}: characterize whether the $10\times$ global-max
gain (\S2.3) yields usable, non-seizing dynamics or merely pushes the failure
threshold higher.
\end{enumerate}

\section{Reproducibility}

All numbers in this note were produced on this repository (Node v24, Python
3.13, Windows; CPU-only; minutes), \textbf{except where a row is explicitly
marked as a historical record}. Commit lineage (v1 $\to$ this revision) is
given on the title line. Vendor data acquisition with pinned commits and
SHA-256 checksums: \texttt{game/tools/fetch\_vendor\_data.md}. Probe batteries
are frozen as scripts (\texttt{probe\_flywire.mjs},
\texttt{probe\_malecns.mjs}, \texttt{probe\_gain.mjs},
\texttt{probe\_signflip.mjs}, \texttt{probe\_supplementary.mjs},
\texttt{reproduce\_sweeps.mjs}); key per-run outputs are archived under
\texttt{game/docs/results/} with the producing command in each file's header.
Timing fields (ms/tick) fluctuate with machine load and are not bit-stable;
all spike counts are.

\begin{center}
\small
\begin{tabular}{L{8.6cm} L{6.2cm}}
\toprule
number(s) & command \\
\midrule
FlyWire propagation (47,106 / 35,159 central); N/E counts & \texttt{node game/test/sim.test.mjs} \\
FlyWire weight distribution; reference-gain failure; targetInput sweep; bias table & \texttt{node game/tools/probe\_flywire.mjs} \\
gain $\times$1/$\times$10 retest, both datasets (16 / 131+13 central) & \texttt{node game/tools/probe\_gain.mjs} \\
FlyWire pools generation & \texttt{python game/tools/prepare\_pools.py} \\
MaleCNS three-file conversion & \texttt{python game/tools/flyb\_to\_bin.py} \\
MaleCNS probes 1--4 & \texttt{node game/tools/probe\_malecns.mjs} \\
MaleCNS GLUT-excitatory variant build & \texttt{python game/tools/flyb\_to\_bin.py {-}{-}glut-excitatory} \\
FlyWire GLUT-inhibitory rebuild + equivalence self-check & \texttt{python game/tools/build\_flywire\_variant.py} \\
sign-flip bidirectional probe & \texttt{node game/tools/probe\_signflip.mjs} \\
MaleCNS feeding/escape/ascending probes & \texttt{node game/tools/probe\_supplementary.mjs} \\
FlyWire closed loop (seed 7: first eat 34.8 s, score 2/120 s) & \texttt{node game/tools/headless\_run.mjs 120 '\{\}' 7 flywire} \\
MaleCNS closed loop (score 1, first eat 104.3 s) & \texttt{node game/tools/headless\_run.mjs 180 '\{\}' 7 malecns} \\
sweep round 1 & \texttt{node game/tools/tune\_sweep.mjs '\{"turnGain":[2.6,4,6],"smellSigma":[130,250],"baseSpeed":[6,12],"bananaMaxDist":[null]\}'} \\
sweep round 2 & \texttt{node game/tools/tune\_sweep.mjs '\{"turnGain":[1.8,2.6],"emaTau":[5,20,40],"baseSpeed":[6,12],"bananaMaxDist":[null]\}'} \\
champion config, historical default restored (3.00; 2.20 under current default) & \texttt{node game/tools/tune\_sweep.mjs '\{"turnGain":[1.8],"emaTau":[5],"baseSpeed":[6],"bananaMaxDist":[null]\}'} \\
out-of-sample \& level-design comparisons (reproduces the \S4.2 table exactly) & \texttt{node game/tools/reproduce\_sweeps.mjs} \\
\bottomrule
\end{tabular}
\end{center}

Notes on exactness. \texttt{"bananaMaxDist":[null]} in the sweep commands
explicitly restores the historical unbounded spawn distance (\texttt{null}
$\to\infty$ in \texttt{harness.mjs}); without it, today's default (250)
silently changes the historical tables --- the reviewer's 2.20-vs-3.00
discrepancy was exactly this. Seeded runs use an injectable LCG replacing
\texttt{Math.random}; episodes call \texttt{sim.reset()} between runs so one
process reuses a single parsed brain. Headless runs are not real-time
throttled. Honest reuse note: seeds 101--110 served both for selecting
\texttt{bananaMaxDist} and for evaluating it (\S4.3); first-eat statistics
count successful episodes only. One row is a \textbf{historical record not
rebuilt by script}: the 350 px no-effect comparison of \S4.2 (its
configuration is documented in the \texttt{game/README.md} tuning log); every
other number in the tables above is reproduced by the listed commands today.

\begin{thebibliography}{99}
\footnotesize
\bibitem{bates2026} Bates, A. S., Phelps, J. S., et al. (2026). Distributed
control circuits across a brain-and-cord connectome. \emph{Nature}.
\url{https://doi.org/10.1038/s41586-026-10735-w}
\bibitem{flybench} brandoncho369/flybench (2026). Behavioural benchmark for
whole-brain fruit fly connectome simulations: pre-registered, cited
behavioural tasks on FlyWire v783 and MaleCNS with shuffled-wiring controls
and a reference LIF implementation.
\url{https://github.com/brandoncho369/flybench} (MIT)
\bibitem{fbm} blendi-remade/fly-brain-minecraft (2026). Minecraft mod driven
by the MaleCNS connectome; source of the FLYB v1 binary packaging of neuPrint
male-cns:v1.0 used here (format spec and full transformation log in its
\texttt{PROVENANCE.md}).
\url{https://github.com/blendi-remade/fly-brain-minecraft} (MIT)
\bibitem{beiran2025} Beiran, M. \& Litwin-Kumar, A. (2025). \emph{Nature
Neuroscience} 28: 2561--2574（连接组只部分约束动力学）. Code:
\url{https://github.com/emebeiran/connconstr} ; data:
\url{https://zenodo.org/records/16618353}
\bibitem{berg2026} Berg, S., Beckett, I., et al. (2026). Sexual dimorphism in
the complete connectome of the Drosophila male central nervous system
(MaleCNS). \emph{Cell}. \url{https://doi.org/10.1016/j.cell.2026.08.015} ;
preprint: \url{https://www.biorxiv.org/content/10.1101/2025.10.09.680999v1}
\bibitem{awesomefly} cobanov/awesome-fly (2026). Community index of fly-brain
demos, with fidelity disclaimers.
\url{https://github.com/cobanov/awesome-fly}
\bibitem{dorkenwald2024} Dorkenwald, S., Matsliah, A., Sterling, A. R., et al.
(2024). Neuronal wiring diagram of an adult brain. \emph{Nature} 634:
124--138. \url{https://doi.org/10.1038/s41586-024-07558-y}
\bibitem{eckstein2024} Eckstein, N., Bates, A. S., Champion, A., et al.
(2024). Neurotransmitter classification from electron microscopy images at
synaptic sites in Drosophila melanogaster. \emph{Cell} 187: 2574--2594.e23.
\url{https://doi.org/10.1016/j.cell.2024.03.016}
\bibitem{schlegel2024} Schlegel, P., Yin, Y., Bates, A. S., et al. (2024).
Whole-brain annotation and multi-connectome cell typing of Drosophila.
\emph{Nature} 634: 139--152.
\url{https://doi.org/10.1038/s41586-024-07686-5}
\bibitem{shiu2024} Shiu, P. K., Sterne, G. R., Spiller, N., et al. (2024). A
Drosophila computational brain model reveals sensorimotor processing.
\emph{Nature} 634: 210--219.
\url{https://doi.org/10.1038/s41586-024-07763-9} ; code:
\url{https://github.com/philshiu/Drosophila_brain_model}
\bibitem{snedea} snedea/flybrain (2026). Browser-based LIF simulation of the
FlyWire connectome (data packaging \& reference simulator; neurotransmitter
sign table in \texttt{scripts/build\_connectome.py}).
\url{https://github.com/snedea/flybrain}
\bibitem{sturner2025} Stürner, T., et al. (2025). \emph{Nature} 643:
158--172（下行/上行神经元的比较连接组学：FlyWire + MANC + FANC）.
\url{https://doi.org/10.1038/s41586-025-08925-z}
\bibitem{wangchen2024} Wang-Chen, S., Stimpfling, A., Lam, K., et al. (2024).
NeuroMechFly v2: simulating embodied sensorimotor control in adult Drosophila.
\emph{Nature Methods} 21: 2353--2362.
\url{https://doi.org/10.1038/s41592-024-02497-y}
\bibitem{zheng2018} Zheng, Z., Lauritzen, J. S., Perlman, E., et al. (2018). A
complete electron microscopy volume of the brain of adult Drosophila
melanogaster. \emph{Cell} 174: 730--743.e22.
\url{https://doi.org/10.1016/j.cell.2018.06.019}
\end{thebibliography}

\end{document}
